\section{5. Conclusion}


In this paper, we propose a novel Unified Dual-Mask Physical Model for robust light field (LF) depth estimation, specifically targeting the formidable challenge posed by ubiquitous non-Lambertian surfaces. Departing from conventional data-driven approaches that heavily rely on the flawed Lambertian assumption, our framework establishes a physically interpretable paradigm by unifying the imaging models of Lambertian, non-Lambertian, and mixed-material scenes. By integrating an MRI-inspired angular frequency analysis mechanism with reverse wavevector parsing, the proposed method effectively disentangles complex reflectance components and adaptively recovers accurate geometry. Experiments validate our algorithm on both standard synthetic benchmarks and real-world LF datasets, demonstrating that our physically grounded approach significantly outperforms state-of-the-art baselines, particularly in regions with extreme specularities and spatially varying BRDFs.

Specifically, the core contributions and technical advancements of this work are summarized as follows:

*   \textbf{MRI-Inspired Angular Frequency Analysis Mechanism:} We introduce a novel frequency-domain analysis technique that overcomes the severe accuracy bottleneck (a mere 24.6\% accuracy) of traditional learning-based material classifiers constrained by small receptive fields (e.g., \(9 \times 9\) patches). Operating with an optimal \(\mathcal{O}(HW)\) computational complexity, this mechanism achieves real-time, high-precision reflectance type identification, providing a reliable physical prior for subsequent depth inference in non-Lambertian scenarios.
*   \textbf{Dual-Mask Physical Modeling and Reverse Wavevector Parsing:} By formulating a dual-mask physical model at the physical optics level, we successfully unify the imaging processes of diverse material properties. The integration of reverse wavevector parsing resolves the inherent lack of physical interpretability in purely data-driven methods. This theoretical guarantee endows our model with exceptional generalization capabilities, enabling it to gracefully handle complex mixed materials and extreme non-Lambertian reflectances without performance degradation.
*   \textbf{Component-Aware Adaptive Depth Estimation Strategy:} To tackle the complete failure of traditional Epipolar Plane Image (EPI) methods—where depth predictions degenerate into texture replication due to the violation of photo-consistency assumptions in non-Lambertian regions—we design a component-aware adaptive strategy. A series of experiments conducted on complex LF datasets validate the superiority of this strategy across multiple dimensions, significantly enhancing depth estimation accuracy and robustness in both mixed and purely non-Lambertian scenes, and yielding substantial reductions in depth error metrics compared to existing methods.

Ultimately, this work not only resolves the long-standing dilemma of non-Lambertian depth estimation in light field imaging but also pioneers a vital shift from purely empirical, data-driven fitting to physically principled modeling. The proposed unified physical framework lays a solid theoretical and practical foundation for advanced 3D vision tasks. The principles established in this study hold immense potential for broader applications in novel view synthesis, high-fidelity 3D reconstruction, and autonomous navigation, paving the way for next-generation dense 3D perception systems that are inherently robust to the complex optical properties of the real world.


From a signal processing perspective, depth estimation in non-Lambertian scenes is fundamentally an ill-posed inverse problem. Traditional photo-consistency metrics fail because specularities introduce high-frequency angular outliers that corrupt the epipolar geometry. Our dual-mask physical model effectively regularizes this problem by decoupling the radiometric and geometric variables. Specifically, the medium mask acts as a spatially varying regularizer that dictates the light-matter interaction regime, while the angular direction mask constrains the wavevector deflection field. By enforcing the Radiance Tensor Field (RTF) rank constraint, defined as:
\begin{align*}
\text{rank}(\mathcal{T}(x, y)) = \begin{cases} 1, \& \text{if } (x, y) \in \Omega_{\text{Lambertian}} \\ \ge 2, \& \text{if } (x, y) \in \Omega_{\text{non-Lambertian}} \end{cases}
\end{align*}
where \(\mathcal{T}(x, y)\) denotes the local radiance tensor, we implicitly embed a structural prior into the optimization landscape. This prevents the depth solution from converging to spurious local minima typically induced by view-dependent highlight shifts.

Despite the robustness of the MRI-inspired angular frequency analysis, the method's efficacy is intrinsically bounded by the angular Nyquist limit. The 2D-DFT operation relies on adequate angular sampling to prevent spectral leakage between the low-frequency diffuse lobe and the high-frequency specular components. While a \(9 \times 9\) angular grid provides sufficient resolution for standard isotropic materials, sparser sampling configurations (e.g., \(5 \times 5\)) inevitably lead to frequency aliasing in the \(k\)-space. In such under-sampled regimes, the mid-frequency scattering components become indistinguishable from broadened specular highlights, degrading the material classification boundary. Furthermore, highly anisotropic materials with narrow, elongated specular lobes introduce directional high-frequency energy that may fall between discrete frequency bins, requiring denser angular sampling to accurately capture the spectral distribution.

Another critical boundary of the proposed physical model lies in its assumption of surface-level Bidirectional Reflectance Distribution Function (BRDF) interactions. The reverse wavevector parsing framework accurately models reflection and surface scattering but struggles with volumetric light transport. For translucent materials exhibiting subsurface scattering, or refractive media such as glass and water, the light rays undergo complex internal bounces and directional bending that violate the surface wavevector deflection assumptions. Consequently, the angular direction mask may yield physically inconsistent deflection vectors for these regions, leading to depth distortions. Extending the current surface-based masks to incorporate volumetric rendering equations remains a necessary direction for handling extreme optical phenomena.

From an engineering and deployment standpoint, while the initial angular frequency analysis achieves an optimal \(\mathcal{O}(HW)\) complexity, the subsequent reverse wavevector parsing introduces non-trivial computational bottlenecks. The constrained least-squares estimation for BRDF parameters and the iterative RTF rank optimization involve dense matrix operations and eigenvalue decompositions that are computationally heavy for high-resolution 4D light fields. To facilitate real-time deployment in latency-critical applications like robotic manipulation, future work must explore algorithmic unrolling or hardware-specific acceleration. Implementing the 2D-DFT and matrix inversion kernels on Field-Programmable Gate Arrays (FPGAs), or mapping the iterative optimization to customized tensor cores, could drastically reduce inference latency, bridging the gap between rigorous physical modeling and edge-computing constraints.

Finally, our findings offer profound implications for the optical design of next-generation light field cameras. The physical model demonstrates that uniform angular sampling is inherently sub-optimal for capturing non-Lambertian scenes, as it wastes spatial resolution on redundant low-frequency diffuse signals while under-sampling critical high-frequency specularities. This insight advocates for the development of adaptive or non-uniform micro-lens arrays that dynamically allocate angular samples based on local scene reflectance, thereby maximizing the information entropy of the captured light field.

Despite the demonstrated efficacy of the unified dual-mask physical model, several limitations remain that outline promising directions for future research. First, the proposed MRI-inspired angular frequency analysis inherently relies on a sufficient angular sampling density (e.g., \(9 \times 9\) sub-aperture views). In scenarios involving sparse light fields (e.g., \(3 \times 3\) or \(5 \times 5\) angular grids), the high-frequency spectral signatures corresponding to specular and scattering components may suffer from severe aliasing, degrading the accuracy of pixel-level material classification. Second, while the reverse wavevector parsing effectively disentangles mixed reflectance, the constrained least-squares BRDF estimation can become ill-posed at extreme depth discontinuities accompanied by severe occlusions. In such regions, the number of valid, unoccluded viewing rays may fall below the theoretical threshold required to robustly resolve the radiation tensor field (RTF) rank and decouple the medium mask from the angular direction mask. Furthermore, the current physical model primarily parameterizes surface reflectance but does not fully account for complex subsurface scattering or translucent materials, where light transport fundamentally violates the localized surface interaction assumption.

To address these challenges, future work will explore integrating implicit neural representations (INRs) to perform physics-aware angular super-resolution prior to the frequency-domain transformation, thereby mitigating spectral aliasing in sparse light fields. Additionally, we aim to extend the dual-mask paradigm into the spatiotemporal domain. By incorporating temporal frequency analysis, the framework could robustly handle dynamic non-Lambertian scenes, leveraging time-coherence to regularize the BRDF parameter estimation in heavily occluded regions. From a systems engineering perspective, we plan to investigate hardware-software co-design strategies, such as deploying customized computational optics or metasurfaces that perform partial angular frequency filtering directly in the optical domain. This would fundamentally offload the backend signal processing and enable true real-time depth perception. Lastly, we intend to explore the fusion of our physical priors with active sensing modalities, such as structured light or time-of-flight (ToF), to provide absolute scale and resolve the inherent scale-depth ambiguity in purely passive light field setups, while incorporating self-supervised photometric losses to enhance zero-shot generalization to in-the-wild captures.